% Chapter 10: Production Pipeline for Causal Inference
% ============================================================================
% This chapter covers:
% - Why production causal inference differs from standard ML
% - Model serialization and versioning for DML
% - Inference APIs for CATE estimation
% - Causal-specific monitoring (overlap, treatment shift, nuisance degradation)
% - Retraining triggers unique to causal inference
% - End-to-end pipeline with Airflow/Prefect
% ============================================================================

\chapter{Production Pipeline for Causal Inference}
\label{ch:production}

\section{Introduction}\label{sec:ch10_introduction}

Chapters 1--9 developed the theory, methodology, and validation of Double Machine Learning for causal inference. This chapter bridges the gap between research and production deployment. Deploying causal inference systems differs fundamentally from standard machine learning in ways that practitioners must understand to build reliable systems.

\keyresult{
\textbf{Chapter Objectives:}
\begin{itemize}
    \item Understand why causal inference deployment differs from standard ML
    \item Design model serialization for DML's two-stage architecture
    \item Build inference APIs for treatment effect estimation
    \item Implement causal-specific monitoring (not just prediction accuracy)
    \item Create intelligent retraining triggers based on causal validity
    \item Assemble end-to-end production pipelines
\end{itemize}
}

\subsection{The Production Paradox}

Standard machine learning production systems monitor prediction accuracy against ground truth. A deployed image classifier can be evaluated against labeled images. A recommendation system can be evaluated by click-through rates. The feedback loop is direct: prediction $\rightarrow$ observation $\rightarrow$ accuracy.

Causal inference breaks this paradigm fundamentally:

\begin{keyinsight}[The Fundamental Problem of Production Causal Inference]
We \textbf{cannot directly monitor treatment effect accuracy} in production because counterfactual outcomes are never observed. A system predicting ``treatment would increase retention by 5\%'' cannot be validated by observing both outcomes for the same unit.
\end{keyinsight}

This creates a production monitoring challenge without parallel in standard ML. Instead of monitoring prediction accuracy, we must monitor the \textit{conditions under which our estimates remain valid}:

\begin{enumerate}
    \item \textbf{Overlap violations} (positivity breakdown)
    \item \textbf{Treatment distribution shift} (assignment mechanism change)
    \item \textbf{Nuisance model degradation} (first-stage fit quality)
    \item \textbf{Effect stability} (is the treatment effect itself changing?)
\end{enumerate}

\begin{table}[htbp]
\centering
\caption{Production monitoring: Standard ML vs. Causal Inference}
\label{tab:ml_vs_causal_monitoring}
\begin{tabular}{lll}
\toprule
\textbf{Aspect} & \textbf{Standard ML} & \textbf{Causal Inference} \\
\midrule
Ground truth & Observable & Never observable \\
Primary metric & Prediction accuracy & Nuisance fit quality \\
Feature drift & Monitor input distribution & Monitor treatment distribution \\
Overlap & N/A & Critical (positivity) \\
Retraining trigger & Accuracy degradation & Identification breakdown \\
Validation & Holdout set & Refutation tests \\
\bottomrule
\end{tabular}
\end{table}

\section{Model Serialization for DML}\label{sec:model_serialization}

DML models have a unique architecture requiring specialized serialization:

\subsection{Components to Serialize}

A production DML model consists of:

\begin{enumerate}
    \item \textbf{Nuisance models} for each cross-fitting fold
    \begin{itemize}
        \item Propensity model: $\hat{e}(X) \approx P(T=1|X)$
        \item Outcome model: $\hat{m}(X) \approx \mathbb{E}[Y|X]$
    \end{itemize}
    \item \textbf{Cross-fitting configuration} (K folds, temporal blocking)
    \item \textbf{Feature transformer} (scaling, encoding)
    \item \textbf{HAC parameters} for time series standard errors
    \item \textbf{Training metadata} (data snapshot, hyperparameters)
\end{enumerate}

% DML model version structure
\begin{minted}[fontsize=\small]{python}
@dataclass
class DMLModelVersion:
    """Versioned DML model with full serialization support."""

    version_id: str  # Unique hash-based identifier
    created_at: str
    model_type: str  # "double_ml", "dynamic_dml", etc.
    n_folds: int

    # Serialized nuisance models by fold
    nuisance_models: Dict[int, Tuple[bytes, bytes]]

    # Feature specification
    feature_names: List[str]
    treatment_name: str
    outcome_name: str

    # Reproducibility
    hyperparameters: Dict[str, Any]
    metrics: Dict[str, float]  # R² propensity, R² outcome

    @classmethod
    def create(
        cls,
        model_type: str,
        nuisance_models: Dict[int, Tuple[Any, Any]],
        feature_names: List[str],
        treatment_name: str,
        outcome_name: str,
        hyperparameters: Dict[str, Any],
        **kwargs,
    ) -> "DMLModelVersion":
        """Create versioned model with content-based hash."""
        # Serialize nuisance models per fold
        serialized = {}
        for fold_idx, (prop, out) in nuisance_models.items():
            serialized[fold_idx] = (
                pickle.dumps(prop),
                pickle.dumps(out),
            )

        # Generate version ID from content hash
        timestamp = datetime.utcnow().isoformat()
        version_id = f"dml-{compute_hash(...)[:12]}"

        return cls(version_id=version_id, ...)
\end{minted}

\subsection{Model Registry Operations}

A DML-aware model registry supports:

\begin{itemize}
    \item \textbf{Registration}: Store versioned models with full metadata
    \item \textbf{Retrieval}: Load specific version for inference
    \item \textbf{Promotion}: Staging $\rightarrow$ Production workflow
    \item \textbf{Rollback}: Revert to previous production version
    \item \textbf{Lineage}: Track which data/config produced which model
\end{itemize}

% Model registry workflow
\begin{minted}[fontsize=\small]{python}
# Initialize registry
registry = DMLModelRegistry("./models/dml_registry")

# After training
version_id = registry.register(model_version)

# Promotion workflow
registry.promote_to_staging(version_id)

# After validation passes
registry.promote_to_production()

# If issues detected
registry.rollback(to_version="dml-abc123")
\end{minted}

\subsection{Persistence Format}

Models are persisted as directories containing:

\begin{verbatim}
dml-abc123xyz/
├── metadata.json          # Version info, hyperparameters
├── propensity_fold_0.pkl  # Serialized propensity model
├── outcome_fold_0.pkl     # Serialized outcome model
├── propensity_fold_1.pkl
├── outcome_fold_1.pkl
├── ...
└── feature_transformer.pkl  # Optional scaler/encoder
\end{verbatim}

This structure enables:
\begin{itemize}
    \item Atomic updates (complete version or nothing)
    \item Efficient fold-level updates (not full model)
    \item Human-readable metadata inspection
    \item Git-compatible versioning (track \texttt{metadata.json})
\end{itemize}


\section{Inference APIs}\label{sec:inference_apis}

Production inference for causal models differs from standard ML prediction.

\subsection{Batch vs. Real-time Inference}

\begin{table}[htbp]
\centering
\caption{Inference patterns for causal models}
\label{tab:inference_patterns}
\begin{tabular}{lll}
\toprule
\textbf{Pattern} & \textbf{Use Case} & \textbf{Latency} \\
\midrule
Batch scoring & Periodic CATE updates for all customers & Minutes \\
Real-time API & On-demand treatment decisions & <100ms \\
Streaming & Continuous effect monitoring & Sub-second \\
\bottomrule
\end{tabular}
\end{table}

\subsection{API Design for CATE Estimation}

A production CATE API should provide:

% CATE estimation API
\begin{minted}[fontsize=\small]{python}
class CATEInferenceAPI:
    """Production API for treatment effect estimation."""

    def estimate_cate(
        self,
        features: Dict[str, Any],
        return_confidence: bool = True,
    ) -> CATEResponse:
        """
        Estimate conditional treatment effect.

        Args:
            features: Covariate values for the unit
            return_confidence: Include confidence intervals

        Returns:
            CATEResponse with effect estimate and metadata
        """
        X = self._transform_features(features)

        # Compute CATE using averaged nuisance models
        cate = self._compute_cate(X)

        # Compute confidence interval
        ci_lower, ci_upper = self._compute_ci(X, cate)

        # Check overlap (propensity extremity)
        propensity = self._predict_propensity(X)
        overlap_warning = self._check_overlap(propensity)

        return CATEResponse(
            cate=cate,
            ci_lower=ci_lower,
            ci_upper=ci_upper,
            propensity=propensity,
            overlap_warning=overlap_warning,
            model_version=self._current_version,
        )
\end{minted}

\subsection{Response Structure}

A complete CATE response includes uncertainty and validity flags:

% CATE response structure
\begin{minted}[fontsize=\small]{python}
@dataclass
class CATEResponse:
    cate: float           # Point estimate
    ci_lower: float       # 95% CI lower
    ci_upper: float       # 95% CI upper
    propensity: float     # P(T=1|X) for this unit
    overlap_warning: bool # True if propensity extreme
    model_version: str    # For reproducibility
    timestamp: str        # Inference timestamp
\end{minted}

\begin{keyinsight}[Always Return Propensity with CATE]
Including the propensity score in CATE responses enables downstream systems to filter or down-weight estimates with poor overlap. A CATE estimate with $\hat{e}(X) = 0.001$ is extrapolation, not estimation.
\end{keyinsight}


\section{Causal-Specific Monitoring}\label{sec:causal_monitoring}

The heart of production causal inference is monitoring the \textit{conditions for valid inference}, not prediction accuracy.

\subsection{Overlap Violations (Positivity)}

The positivity assumption requires $0 < P(T=1|X) < 1$ for all $X$. In practice, we need propensity scores bounded away from 0 and 1:

\begin{equation}
\epsilon < \hat{e}(X) < 1 - \epsilon, \quad \text{typically } \epsilon = 0.01
\end{equation}

% Overlap monitoring
\begin{minted}[fontsize=\small]{python}
def check_overlap_violations(
    propensity_scores: np.ndarray,
    clip_min: float = 0.01,
    clip_max: float = 0.99,
) -> MonitoringResult:
    """Check for positivity violations."""

    too_low = propensity_scores < clip_min
    too_high = propensity_scores > clip_max
    violation_rate = (np.sum(too_low) + np.sum(too_high)) / len(propensity_scores)

    if violation_rate >= 0.10:  # >10% violations
        level = AlertLevel.CRITICAL
        message = f"CRITICAL: {violation_rate:.1%} overlap violations"
    elif violation_rate >= 0.05:
        level = AlertLevel.WARNING
        message = f"WARNING: {violation_rate:.1%} overlap violations"
    else:
        level = AlertLevel.OK
        message = f"Overlap OK: {violation_rate:.1%} violations"

    return MonitoringResult(
        check_name="overlap_violations",
        level=level,
        value=violation_rate,
        message=message,
    )
\end{minted}

\subsection{Treatment Distribution Shift}

Unlike standard feature drift, treatment distribution shift indicates the \textit{assignment mechanism} may have changed. This is causal-specific:

\begin{itemize}
    \item Standard ML: Features changed $\rightarrow$ model may underfit
    \item Causal inference: Treatment assignment changed $\rightarrow$ \textbf{identification may fail}
\end{itemize}

% Treatment shift monitoring
\begin{minted}[fontsize=\small]{python}
def check_treatment_shift(
    treatment_current: np.ndarray,
    treatment_baseline: np.ndarray,
) -> MonitoringResult:
    """Detect treatment distribution shift."""

    # For binary treatment: proportion difference
    # For continuous: KS statistic
    if is_binary(treatment_current):
        p_current = np.mean(treatment_current)
        p_baseline = np.mean(treatment_baseline)
        metric = abs(p_current - p_baseline)
    else:
        metric, _ = ks_2samp(treatment_current, treatment_baseline)

    if metric >= 0.10:
        level = AlertLevel.CRITICAL
        message = "Treatment assignment mechanism may have changed"
    # ...
\end{minted}

\subsection{Nuisance Model Degradation}

DML validity requires well-fitting nuisance models. We monitor nuisance $R^2$, not treatment effect accuracy (which is unknowable):

% Nuisance degradation monitoring
\begin{minted}[fontsize=\small]{python}
def check_nuisance_degradation(
    r2_propensity: float,
    r2_outcome: float,
    warning_threshold: float = 0.50,
    critical_threshold: float = 0.30,
) -> MonitoringResult:
    """Check nuisance model fit quality."""

    min_r2 = min(r2_propensity, r2_outcome)

    if min_r2 < critical_threshold:
        level = AlertLevel.CRITICAL
        message = f"Nuisance fit poor (R²={min_r2:.2f}). DML may be biased."
    elif min_r2 < warning_threshold:
        level = AlertLevel.WARNING
        message = f"Nuisance fit degraded (R²={min_r2:.2f})"
    else:
        level = AlertLevel.OK

    return MonitoringResult(...)
\end{minted}

\begin{keyinsight}[Why We Monitor Nuisance $R^2$ Instead of Effect Accuracy]
The DML bias term depends on the product of nuisance errors:
\[
\text{Bias} \propto \|\hat{m} - m_0\| \cdot \|\hat{e} - e_0\|
\]
Poor nuisance fit indicates potential bias. But we cannot directly measure treatment effect accuracy because counterfactuals are unobserved.
\end{keyinsight}

\subsection{Effect Stability}

Large changes in estimated treatment effects over time warrant investigation:

% Effect stability monitoring
\begin{minted}[fontsize=\small]{python}
def check_effect_stability(
    current_effect: float,
    baseline_effect: float,
    current_se: Optional[float] = None,
    baseline_se: Optional[float] = None,
) -> MonitoringResult:
    """Monitor treatment effect stability over time."""

    if abs(baseline_effect) < 1e-10:
        relative_change = float("inf") if abs(current_effect) > 1e-10 else 0.0
    else:
        relative_change = abs(current_effect - baseline_effect) / abs(baseline_effect)

    if relative_change >= 0.50:  # 50% change
        level = AlertLevel.CRITICAL
        message = f"Effect changed {relative_change:.0%}: investigate cause"
    elif relative_change >= 0.20:  # 20% change
        level = AlertLevel.WARNING
    # ...
\end{minted}

Effect changes may indicate:
\begin{itemize}
    \item Genuine time-varying effects (valid, requires dynamic modeling)
    \item Model degradation (problematic)
    \item Selection/survivorship bias in new data (problematic)
    \item Treatment definition changed (problematic)
\end{itemize}


\section{Retraining Triggers}\label{sec:retraining_triggers}

Standard ML retraining triggers focus on prediction accuracy degradation. Causal inference requires different triggers aligned with identification assumptions.

\subsection{Causal-Specific Triggers}

\begin{table}[htbp]
\centering
\caption{Retraining triggers: Standard ML vs. Causal Inference}
\label{tab:retrain_triggers}
\begin{tabular}{lll}
\toprule
\textbf{Standard ML} & \textbf{Causal Inference} & \textbf{Rationale} \\
\midrule
Accuracy drop & N/A & Counterfactuals unobserved \\
Feature drift & Covariate drift & Affects nuisance models \\
Class imbalance & Treatment shift & Assignment mechanism changed \\
--- & Overlap violations & Positivity breakdown \\
--- & Nuisance $R^2$ drop & DML validity threatened \\
Scheduled & Scheduled & Preventive maintenance \\
\bottomrule
\end{tabular}
\end{table}

\subsection{Scheduler Design}

% Retraining scheduler
\begin{minted}[fontsize=\small]{python}
class RetrainScheduler:
    """Intelligent retraining scheduler for DML pipelines."""

    def __init__(self, config: RetrainSchedulerConfig):
        self.config = config
        self.monitor = CausalMonitor()

    def evaluate_retrain_need(
        self,
        monitoring_results: List[MonitoringResult],
        n_samples: int,
    ) -> Optional[RetrainTrigger]:
        """Evaluate if retraining is needed."""

        # Check cooldown (avoid excessive retraining)
        if self.is_in_cooldown():
            return None

        # Check minimum sample threshold
        if n_samples < self.config.min_samples:
            return None

        # Map monitoring results to triggers
        for result in monitoring_results:
            if result.level == AlertLevel.CRITICAL:
                return RetrainTrigger(
                    trigger_type=self._result_to_trigger_type(result),
                    severity=AlertLevel.CRITICAL,
                    reason=result.message,
                )

        # Check scheduled retraining
        return self.check_scheduled_retrain()
\end{minted}

\subsection{Trigger Types}

% Trigger type enumeration
\begin{minted}[fontsize=\small]{python}
class TriggerType(Enum):
    SCHEDULED = "scheduled"           # Time-based
    OVERLAP_VIOLATION = "overlap"     # Positivity breakdown
    TREATMENT_SHIFT = "treatment"     # Assignment mechanism
    NUISANCE_DEGRADATION = "nuisance" # First-stage fit
    EFFECT_INSTABILITY = "effect"     # Effect magnitude change
    COVARIATE_SHIFT = "covariate"     # Feature distribution
    MANUAL = "manual"                 # Human-initiated
\end{minted}


\section{End-to-End Pipeline}\label{sec:production_pipeline}

This section assembles all components into a production-ready pipeline.

\subsection{Pipeline Architecture}

\begin{figure}[htbp]
\centering
\begin{tikzpicture}[
    node distance=2cm,
    box/.style={rectangle, draw, minimum width=3cm, minimum height=1cm, align=center},
    decision/.style={diamond, draw, minimum width=2cm, minimum height=1cm, align=center, aspect=2}
]

% Nodes
\node[box] (data) {Data Ingestion\\(FRED + Insurance)};
\node[box, below of=data] (features) {Feature Engineering};
\node[box, below of=features] (dml) {DML Estimation\\(Cross-fitting)};
\node[box, below of=dml] (monitor) {Causal Monitoring};
\node[decision, below of=monitor] (trigger) {Retrain?};
\node[box, right of=trigger, xshift=3cm] (retrain) {Retrain Pipeline};
\node[box, below of=trigger, yshift=-1cm] (serve) {Model Serving};

% Edges
\draw[->] (data) -- (features);
\draw[->] (features) -- (dml);
\draw[->] (dml) -- (monitor);
\draw[->] (monitor) -- (trigger);
\draw[->] (trigger) -- node[above] {yes} (retrain);
\draw[->] (trigger) -- node[right] {no} (serve);
\draw[->] (retrain) |- (dml);

\end{tikzpicture}
\caption{Production DML pipeline architecture}
\label{fig:pipeline_architecture}
\end{figure}

\subsection{InsuranceDMLPipeline Implementation}

% Complete production pipeline
\begin{minted}[fontsize=\small]{python}
class InsuranceDMLPipeline:
    """Production DML pipeline for insurance pricing analysis."""

    def __init__(self, config: PipelineConfig):
        self.config = config
        self._registry = DMLModelRegistry(config.model_registry_path)
        self._monitor = CausalMonitor()
        self._scheduler = RetrainScheduler(monitor=self._monitor)

    def fit(
        self,
        X: np.ndarray,
        T: np.ndarray,
        Y: np.ndarray,
        baseline_ate: Optional[float] = None,
    ) -> PipelineResult:
        """
        Fit DML pipeline to data.

        Stages:
        1. Data preparation (scaling)
        2. Cross-fitted nuisance estimation
        3. ATE estimation with orthogonalization
        4. HAC standard errors (time series)
        5. Monitoring checks
        6. Model versioning
        """
        # Prepare data
        X_scaled, T, Y = self._prepare_data(X, T, Y)

        # Create cross-validator
        cv = self._create_cross_validator(len(T))

        # Fit nuisance models per fold
        propensity_scores = np.zeros(len(T))
        outcome_predictions = np.zeros(len(T))
        nuisance_models = {}

        for fold_idx, (train_idx, test_idx) in enumerate(cv.split(X_scaled)):
            prop_model, out_model = self._fit_fold(
                X_scaled, T, Y, train_idx, test_idx
            )
            propensity_scores[test_idx] = self._predict_propensity(
                prop_model, X_scaled[test_idx]
            )
            outcome_predictions[test_idx] = out_model.predict(X_scaled[test_idx])
            nuisance_models[fold_idx] = (prop_model, out_model)

        # Compute ATE via orthogonalization
        ate, ate_se = self._compute_ate(
            Y, outcome_predictions, T, propensity_scores
        )

        # Run monitoring
        monitoring_results = self._monitor.run_all_checks(
            propensity_scores=propensity_scores,
            r2_propensity=self._compute_r2(propensity_scores, T),
            r2_outcome=self._compute_r2(outcome_predictions, Y),
            current_effect=ate,
            baseline_effect=baseline_ate,
        )

        # Create and register model version
        version = DMLModelVersion.create(
            model_type="double_ml",
            nuisance_models=nuisance_models,
            feature_names=self.config.feature_columns,
            treatment_name=self.config.treatment_column,
            outcome_name=self.config.outcome_column,
            hyperparameters=self._get_hyperparameters(),
        )
        version_id = self._registry.register(version)

        return PipelineResult(
            ate=ate,
            ate_se=ate_se,
            ate_ci_lower=ate - 1.96 * ate_se,
            ate_ci_upper=ate + 1.96 * ate_se,
            nuisance_metrics={"r2_propensity": ..., "r2_outcome": ...},
            monitoring_results=monitoring_results,
            model_version=version_id,
        )
\end{minted}

\subsection{Workflow Orchestration}

For production, integrate with Airflow or Prefect:

% Airflow DAG skeleton
\begin{minted}[fontsize=\small]{python}
# airflow_dags/dml_retrain.py

from airflow import DAG
from airflow.operators.python import PythonOperator, BranchPythonOperator
from datetime import datetime, timedelta

default_args = {
    "owner": "causal-ml",
    "retries": 1,
    "retry_delay": timedelta(minutes=5),
}

with DAG(
    dag_id="dml_retrain_pipeline",
    default_args=default_args,
    schedule_interval="@daily",
    start_date=datetime(2024, 1, 1),
    catchup=False,
) as dag:

    load_data = PythonOperator(
        task_id="load_data",
        python_callable=load_inference_data,
    )

    run_monitoring = PythonOperator(
        task_id="run_monitoring",
        python_callable=run_causal_monitoring,
    )

    evaluate_retrain = BranchPythonOperator(
        task_id="evaluate_retrain",
        python_callable=evaluate_retrain_trigger,
    )

    retrain_model = PythonOperator(
        task_id="retrain_model",
        python_callable=retrain_dml_pipeline,
    )

    validate_model = PythonOperator(
        task_id="validate_model",
        python_callable=validate_new_model,
    )

    promote_staging = PythonOperator(
        task_id="promote_staging",
        python_callable=promote_to_staging,
    )

    skip_retrain = PythonOperator(
        task_id="skip_retrain",
        python_callable=lambda: None,
    )

    # Define dependencies
    load_data >> run_monitoring >> evaluate_retrain
    evaluate_retrain >> [retrain_model, skip_retrain]
    retrain_model >> validate_model >> promote_staging
\end{minted}


\section{Insurance Pricing Case Study}\label{sec:insurance_case_study}

We conclude with a complete example applying the production pipeline to insurance competitor pricing.

\subsection{Problem Setup}

An insurance company wants to estimate how competitor pricing affects policy retention:
\begin{itemize}
    \item \textbf{Treatment}: Competitor price difference (continuous)
    \item \textbf{Outcome}: Retention indicator (binary)
    \item \textbf{Confounders}: Customer demographics, policy features, macroeconomic conditions
    \item \textbf{Time series}: Quarterly observations with autocorrelation
\end{itemize}

\subsection{Pipeline Configuration}

% Insurance pipeline configuration
\begin{minted}[fontsize=\small]{python}
from sklearn.ensemble import GradientBoostingClassifier, GradientBoostingRegressor
from src.production import InsuranceDMLPipeline, PipelineConfig

config = PipelineConfig(
    n_folds=5,
    model_registry_path="./models/insurance_dml",

    # Nuisance models
    propensity_model=GradientBoostingClassifier(
        n_estimators=100,
        max_depth=4,
        random_state=42,
    ),
    outcome_model=GradientBoostingRegressor(
        n_estimators=100,
        max_depth=4,
        random_state=42,
    ),

    # Time series configuration
    use_hac=True,
    time_column="quarter",

    # Feature specification
    feature_columns=[
        "age", "income", "tenure", "coverage_amount",
        "gdp_growth", "unemployment", "interest_rate",
    ],
    treatment_column="competitor_price_diff",
    outcome_column="retained",
)

pipeline = InsuranceDMLPipeline(config)
\end{minted}

\subsection{Training and Monitoring}

% Training with monitoring
\begin{minted}[fontsize=\small]{python}
# Load data
X, T, Y = load_insurance_data()

# Fit pipeline
result = pipeline.fit(X, T, Y)

print(f"ATE: {result.ate:.4f}")
print(f"95% CI: [{result.ate_ci_lower:.4f}, {result.ate_ci_upper:.4f}]")
print(f"Nuisance R² (propensity): {result.nuisance_metrics['r2_propensity']:.3f}")
print(f"Nuisance R² (outcome): {result.nuisance_metrics['r2_outcome']:.3f}")

# Check monitoring
for check in result.monitoring_results:
    if not check.is_ok():
        print(f"ALERT: {check.message}")

# Promote to production
pipeline.promote_to_production()
\end{minted}

\subsection{Production Inference}

% Production inference
\begin{minted}[fontsize=\small]{python}
# Evaluate retraining need with new data
X_new, T_new = load_new_quarter_data()

results, trigger = pipeline.evaluate_retrain_need(
    X_new, T_new,
    X_baseline=X[:len(X_new)],
    T_baseline=T[:len(X_new)],
)

if trigger is not None:
    print(f"Retrain triggered: {trigger.reason}")
    # Initiate retraining workflow
else:
    print("Model stable - no retraining needed")

    # Predict propensity for new customers
    propensity = pipeline.predict_propensity(X_new)

    # Flag high-risk customers (extreme propensity)
    high_risk = (propensity < 0.05) | (propensity > 0.95)
    print(f"High-risk customers: {np.sum(high_risk)} ({np.mean(high_risk):.1%})")
\end{minted}


\section{Exercises}\label{sec:ch10_exercises}

\begin{exercise}[Monitoring Dashboard]
Design a monitoring dashboard for a production DML system. What visualizations would you include? Consider:
\begin{enumerate}
    \item Propensity score distribution over time
    \item Treatment distribution shift tracking
    \item Nuisance model $R^2$ trends
    \item Effect estimate stability with confidence bands
\end{enumerate}
\end{exercise}

\begin{exercise}[Retraining Threshold Calibration]
The default overlap violation threshold of 10\% for CRITICAL alerts may be too aggressive or lenient depending on the application. Describe a procedure for calibrating this threshold using:
\begin{enumerate}
    \item Monte Carlo simulation with known treatment effects
    \item Historical data from past model deployments
\end{enumerate}
\end{exercise}

\begin{exercise}[A/B Test Integration]
Propose how to integrate randomized A/B tests with a production DML system. How can occasional randomization provide:
\begin{enumerate}
    \item Validation of DML estimates
    \item Ground truth for effect stability monitoring
    \item Retraining triggers when observational estimates diverge from experimental
\end{enumerate}
\end{exercise}

\begin{exercise}[Multi-Treatment Extension]
Extend the production pipeline to handle multiple treatments simultaneously. What additional monitoring is needed when:
\begin{enumerate}
    \item Treatments are mutually exclusive (one of K)
    \item Treatments can be combined (2\textsuperscript{K} combinations)
    \item Treatment assignment is continuous multivariate
\end{enumerate}
\end{exercise}


\section{Summary}\label{sec:ch10_summary}

This chapter developed production infrastructure for causal inference systems:

\keyresult{
\textbf{Key Takeaways:}
\begin{itemize}
    \item Production causal inference \textbf{cannot monitor treatment effect accuracy} directly
    \item Monitor \textbf{identification conditions}: overlap, treatment shift, nuisance fit
    \item DML models require \textbf{specialized serialization} for cross-fitting structure
    \item Retraining triggers must be \textbf{causal-specific}, not accuracy-based
    \item End-to-end pipelines integrate monitoring, versioning, and orchestration
\end{itemize}
}

The production module developed here provides:
\begin{itemize}
    \item \texttt{DMLModelVersion}, \texttt{DMLModelRegistry}: Model versioning
    \item \texttt{CausalMonitor}: Overlap, treatment shift, nuisance, effect monitoring
    \item \texttt{RetrainScheduler}: Causal-specific retraining triggers
    \item \texttt{InsuranceDMLPipeline}: End-to-end integration
\end{itemize}

\vspace{1em}

\begin{tcolorbox}[colback=gray!5, colframe=gray!75, title={Chapter Checklist}]
\begin{itemize}
    \item[$\square$] Understand why causal production differs from standard ML
    \item[$\square$] Design model serialization for DML architecture
    \item[$\square$] Implement causal-specific monitoring (4 checks)
    \item[$\square$] Configure retraining triggers based on identification
    \item[$\square$] Assemble end-to-end pipeline with orchestration
\end{itemize}
\end{tcolorbox}

The next and final appendix provides a roadmap for implementing this book's methodology in Julia, offering performance advantages for large-scale causal inference.
